\documentclass{article}
\usepackage[preprint,nonatbib]{neurips_2025}
\usepackage[T1]{fontenc}
\usepackage[utf8]{inputenc}
\usepackage{amsmath,amssymb,booktabs,graphicx,caption,url,hyperref}
\hypersetup{colorlinks=true,linkcolor=blue,citecolor=blue,urlcolor=blue}
\title{Daily Temperature Forecasting\\with Seasonal Regression and ARMA Errors}
\author{Jiayu \qquad Laman\\CPTS 570 Machine Learning, Fall 2026}
\begin{document}
\maketitle
\begin{abstract}
We compare seasonal regression with ARMA errors against persistence, independent-
error seasonal regression, and nonseasonal ARIMA for daily temperature at
Pullman/Moscow Regional Airport. NOAA GSOD records from 2019--2024 provide
2,191 observed days. We train on 2019--2022, select using 2023 validation,
then refit through 2023 and freeze parameters during rolling 2024 testing.
The selected three-harmonic ARMA(1,1) model has test RMSE of 2.405\,$^\circ$C
at one day and 4.019\,$^\circ$C at seven days. The latter improves on seasonal
regression by only 0.015\,$^\circ$C. Its 95\% interval coverage is 93.7\% and
95.6\%, with weaker winter coverage. Results favor short-term residual modeling
at one day and reveal limits of constant-variance Gaussian intervals.
\end{abstract}

\section{Research problem and objectives}
We forecast daily mean temperature at one eastern Washington station, one and
seven days ahead, using only station history and known future calendar terms.
Annual seasonality and recent departures may provide complementary information. We compare their combination with persistence and each component alone.
Our contribution is a reproducible station-level evaluation of an established method.

\section*{Literature review}
\textbf{Seasonality and correlated errors.} Hyndman and Athanasopoulos [2]
describe dynamic harmonic regression, combining Fourier predictors for long
seasonal periods with ARMA errors for short-term dependence. A small harmonic
count produces a smooth annual curve, while higher harmonics allow more detail.
This motivates our annual basis and bounded complexity search. Its fixed
seasonal shape is also a limitation: annual terms do not anticipate unusual
weather departures. Our component comparisons test whether residual dynamics
add useful information at the two forecast horizons.

\textbf{Automatic ARIMA forecasting.} Hyndman and Khandakar [4] present
automatic univariate forecasting algorithms, including a stepwise ARIMA search.
Their work motivates systematic comparison of model orders rather than an
arbitrary single order. We use a small exhaustive grid and chronological
validation RMSE instead of reproducing their automatic selection procedure.
This makes the selected complexity auditable and keeps the evaluation aligned
with our two horizons.

\textbf{Interpretable additive alternatives.} Taylor and Letham [5] propose
an interpretable modular regression approach for forecasting at scale, with
trend, seasonal, and holiday components. It shares the aim of separating
predictable structure from remaining variation. Our study instead focuses
on explicit ARMA residual dependence and the proposal's component baselines.
We do not fit Prophet or infer superiority over that broader model family.

\newpage
\section{Data and preparation}
\subsection{Dataset and station audit}
We use NOAA's Global Surface Summary of the Day (GSOD) archive [1], retaining
records dated January 1, 2019 through December 31, 2024. GSOD daily summaries
are based on subdaily observations. The target is the reported daily mean
temperature, not the midpoint of daily minimum and maximum temperatures.
NOAA documentation defines the daily summaries using Greenwich Mean Time,
so the forecast calendar here consists of UTC days.

The primary station is Pullman/Moscow Regional Airport: USAF 727857, WBAN
94129, combined GSOD identifier 72785794129, and ICAO KPUW. The station
inventory gives latitude 46.744\,$^\circ$N, longitude 117.116\,$^\circ$W, and
elevation 775.7\,m. We also audited Spokane International Airport
(72785024157), the proposal's backup. Pullman has 99.954\% temperature coverage,
which is adequate to retain the preferred station. Spokane has complete coverage,
but is not used in model comparisons. Station choice is based on availability,
not test-year forecast performance.

\begin{center}
\captionof{table}{Chronological partitions at the selected Pullman station.}
\begin{tabular}{lrrl}
\toprule
Partition & Calendar days & Observed & Dates\\
\midrule
Training & 1,461 & 1,460 & 2019--2022\\
Validation & 365 & 365 & 2023\\
Test & 366 & 366 & 2024\\
\bottomrule
\end{tabular}
\end{center}

\subsection{Cleaning and calendar preservation}
We download station-year CSV files and keep the originals unchanged. Each input
has its source URL and SHA-256 checksum recorded. The preparation script checks
station identity and duplicate dates, removes exact duplicates if present, and
stops on conflicting duplicate observations. There are no duplicate dates in
the selected station files. The temperature missing-value sentinel is 9999.9;
no such values appear in the retrieved Pullman records. The
\texttt{TEMP\_ATTRIBUTES} column is an observation count, not a quality flag.
No supplied Pullman mean has a nonpositive count. Other source attributes remain
available in the raw files.

Temperatures are converted by $C=(F-32)5/9$. We reindex onto the full daily
calendar, preserving leap days and the single absent observation on December 31,
2021. We do not interpolate, smooth, or carry forward that value. State-space
models can propagate the state across an unobserved measurement; independent
seasonal regression uses the observed values. Missing forecast targets would
be excluded from scoring. Both validation and test targets are complete.

\begin{center}
\includegraphics[width=\linewidth]{figures/data.pdf}
\captionof{figure}{Daily temperatures and the fixed evaluation partitions.
The strong annual cycle motivates known future calendar predictors.}
\end{center}

\newpage
\section{Forecasting models}
\subsection{Seasonal regression with ARMA/ARIMA errors}
Let $t$ be elapsed calendar days since January 1, 2019. For a fixed annual
period $P=365.25$, the model is
\begin{align}
y_t &= s_t+u_t,\\
s_t &= \beta_0+\sum_{k=1}^{K}\left[
a_k\sin\left(\frac{2\pi kt}{P}\right)+
b_k\cos\left(\frac{2\pi kt}{P}\right)\right].\label{eq:season}
\end{align}
Each pair adds an interpretable harmonic with period $P/k$.
The fixed fractional period approximates a year; it is not an exact calendar
anniversary transformation. Computing features from elapsed dates avoids
resetting phase at split boundaries or compressing time at missing observations.
The seasonal coefficients remain fixed during evaluation.

The residual process satisfies
\begin{equation}
\phi(B)(1-B)^d u_t=\theta(B)\varepsilon_t,
\quad \varepsilon_t\sim\mathcal{N}(0,\sigma^2),
\end{equation}
where $B$ is the lag operator,
$\phi(B)=1-\sum_{j=1}^{p}\phi_jB^j$, and
$\theta(B)=1+\sum_{j=1}^{q}\theta_jB^j$.
For $d=0$, these are ARMA$(p,q)$ errors. If residual nonstationarity justified
$d=1$, the redundant regression intercept would be omitted. This follows the
proposal's dynamic harmonic regression formulation [2]. In the implemented
combined-model search, training stationarity checks support $d=0$.

We estimate regression and error parameters jointly by maximum likelihood using
\texttt{statsmodels} SARIMAX [3]. The Fourier matrix includes its own intercept,
so the combined model uses no additional SARIMAX trend. Stationarity and
invertibility constraints are enabled. For a future date, Fourier features are
known, while the residual mean is forecast recursively. Thus
\begin{equation}
\widehat y_{t+h\mid t}=\widehat s_{t+h}+\widehat u_{t+h\mid t}.
\end{equation}
Prediction intervals come from the state-space forecast distribution conditional
on the estimated parameters. They represent uncertainty in a future temperature
observation, rather than a confidence interval for the seasonal mean alone.

\subsection{Component comparisons}
\textbf{Persistence} predicts $\widehat y_{t+h\mid t}=y_t$ for either horizon.
It measures how much a model improves over repeating the latest observed
temperature. All evaluated origins have observed temperatures.

\textbf{Seasonal regression} uses Equation~\ref{eq:season} with independent
Gaussian errors, implemented as ARIMA$(0,0,0)$ errors. It captures annual
structure but cannot use recent departures from the seasonal mean.

\textbf{Nonseasonal ARIMA} uses no Fourier terms. We search $p,q=0,1,2,3$
and $d=0,1$. A constant is included when $d=0$ and omitted when $d=1$;
there is no drift in the differenced candidates. This comparison tests whether
short-term dependence without explicit annual terms is sufficient.

The \textbf{combined model} uses both components. Its search includes
ARMA$(0,0)$, deliberately nesting the independent-error seasonal model.
This overlap allows the comparison to reveal whether residual dependence is
useful, rather than forcing it into the selected specification.

\subsection{Selected model after the final refit}
The selected $K=3$ model refitted through 2023 has seasonal intercept
9.254\,$^\circ$C, $\widehat\phi_1=0.621$, $\widehat\theta_1=0.358$, and
$\widehat\sigma^2=5.623$\,$^\circ$C$^2$. For ARMA(1,1), residual mean forecasts
beyond the first step decay by successive powers of $\widehat\phi_1$.
Recent residual information therefore has much less influence at seven days.
All fitted coefficients are saved in the reproducibility artifacts and held
fixed throughout the test.

\newpage
\section{Experimental protocol}
\subsection{Chronological validation and model selection}
The initial fit uses 2019--2022. We assess candidate forecasts at daily origins
through 2023, keeping fitted parameters fixed while updating states. The first
origin is December 31, 2022. A forecast is scored only if its target falls in
2023. Consequently there are 365 one-day and 359 seven-day forecasts; the first
seven-day target is January 7. All models use these same origin--target pairs.

We search $K\in\{1,2,3\}$ and $p,q\in\{0,1,2,3\}$ for combined models.
Training ordinary seasonal-regression residuals are checked with augmented
Dickey--Fuller and KPSS tests. We consider $d=1$ for a harmonic count only if
ADF does not reject a unit root at 5\% and KPSS rejects level stationarity at
5\%. No harmonic count satisfies both conditions. This results in 48 combined
candidates, three independent-error seasonal candidates, and 32 nonseasonal
ARIMA candidates: 83 fits in total. All converge. These stationarity checks
are preliminary diagnostics; their residual input drops the single missing
training value and therefore closes one calendar gap. The likelihood fits and
forecast calendars preserve that gap.

The proposal leaves cross-horizon model selection unspecified. We choose one
specification per family by minimizing the equal-weight mean of the two
validation RMSE values:
\begin{equation}
S(m)=\tfrac12\left(\operatorname{RMSE}_{1}(m)+
\operatorname{RMSE}_{7}(m)\right).
\end{equation}
Both terms use Celsius, so no scale normalization is required. This choice
avoids a different model for every horizon while making the tradeoff explicit.
We record AIC, convergence warnings, parameters, and residual checks, but AIC
does not determine the selected family specifications. Nonconverged fits would
be excluded. The optimizer uses up to 300 iterations with a Powell-assisted
retry if needed, followed by another likelihood optimization.

\subsection{Frozen test and information availability}
After validation, each selected model is refitted on 2019--2023. Model orders
are saved before test evaluation. Parameters are frozen for 2024; state updates
do not change regression or ARMA coefficients. At each origin, the evaluator:
\begin{enumerate}
\item forms known future Fourier terms and predicts up to seven days ahead;
\item saves the exact one-day and seven-day predictions and intervals;
\item reveals the next daily temperature and filters the state forward;
\item advances one calendar day and repeats, without parameter refitting.
\end{enumerate}
The test begins at the December 31, 2023 origin and contains 366 one-day and
360 seven-day forecasts. No intervening observation enters an earlier seven-day
forecast. We assume completed daily summaries are available after each UTC day;
actual data-publication delays are not represented.

\subsection{Metrics and verification}
For errors $e_i=\widehat y_i-y_i$ over $n$ eligible targets, we report
$\operatorname{RMSE}=\sqrt{n^{-1}\sum_i e_i^2}$ and
$\operatorname{MAE}=n^{-1}\sum_i|e_i|$ separately by horizon.
For the combined model, coverage is $n^{-1}\sum_i\mathbf{1}\{L_i\leq y_i\leq U_i\}$
and mean width is $n^{-1}\sum_i(U_i-L_i)$ for nominal 95\% intervals.
Protocol checks verify date alignment and common targets, and confirm that
altering future values does not affect forecasts already issued.

\newpage
\section{Validation results and selected specifications}
Table~\ref{tab:validation} reports the selected model from each family. Seasonal
regression selects three Fourier pairs. The nonseasonal baseline selects
ARIMA$(3,0,3)$, while the combined model selects three pairs with
ARMA$(1,1)$ errors. The latter has seven seasonal coefficients, two error
coefficients, and an innovation variance. The search therefore selects a compact
residual model instead of the maximum permitted ARMA orders.

\begin{center}
\captionof{table}{2023 validation metrics. Each model is fitted on 2019--2022.
One-day metrics use 365 targets and seven-day metrics use 359.}
\label{tab:validation}
\begin{tabular}{lrrrr}
\toprule
& \multicolumn{2}{c}{One day} & \multicolumn{2}{c}{Seven days}\\
Model & RMSE & MAE & RMSE & MAE\\
\midrule
Persistence & 2.441 & 1.881 & 5.248 & 4.039\\
Seasonal ($K=3$) & 3.661 & 2.896 & \textbf{3.672} & \textbf{2.903}\\
ARIMA$(3,0,3)$ & 2.259 & 1.763 & 4.457 & 3.429\\
Combined ($K=3$, ARMA$(1,1)$) & \textbf{2.169} & \textbf{1.703} & 3.677 & 2.906\\
\bottomrule
\end{tabular}
\end{center}

At one day, persistence already outperforms the calendar-only forecast,
indicating that recent conditions contain useful information. Nonseasonal ARIMA
improves on persistence, and adding explicit seasonality improves further. The
combined model's one-day RMSE is 2.169\,$^\circ$C. At seven days, however,
seasonal regression has marginally lower RMSE than the combined model:
3.672 versus 3.677\,$^\circ$C. Seasonal structure is more useful at this horizon
than repeating current temperature or relying on the nonseasonal dynamics.

The equal-weight cross-horizon criterion selects the combined specification
within its family, but the table does not imply that residual dynamics improve
every horizon. The two seasonal models are practically very close at seven
days. This distinction is central to the research question: the evidence for
combining components is strongest at the short horizon.

Training ADF p-values for the three seasonal residual series range from about
$1.2\times10^{-14}$ to $4.5\times10^{-12}$. KPSS reports 0.1 for each;
this is a lookup-table bound, not an exact p-value. These checks support the
decision to omit differencing in the combined-model search.

\begin{center}
\includegraphics[width=\linewidth]{figures/acf.pdf}
\captionof{figure}{ACF and PACF of the selected combined model's standardized
training innovations. The longest continuous observed segment is used after
a 30-day burn-in. Shading is the plotting routine's nominal pointwise band.}
\label{fig:acf}
\end{center}

For that continuous segment, Ljung--Box p-values are 0.517 at lag 14 and 0.169
at lag 30, adjusting degrees of freedom by $p+q$. Remaining correlations are
small in Figure~\ref{fig:acf}, although pointwise bands are not simultaneous
tests over all plotted lags. These results motivate proceeding with the selected
model; they do not establish Gaussianity or constant variance.

\newpage
\section{Held-out point forecast performance}
Table~\ref{tab:test} gives the primary 2024 test results after the one-time
2019--2023 refit. The combined model has the lowest observed RMSE and MAE at
both horizons. Its one-day RMSE is 2.405\,$^\circ$C, compared with 2.578 for
persistence, 4.003 for seasonal regression, and 2.486 for nonseasonal ARIMA.
The combined model thus retains the validation pattern of benefiting from
both the annual mean and recent residual behavior at one day.

\begin{center}
\captionof{table}{2024 held-out point forecast metrics in Celsius. The common
sample has 366 one-day and 360 seven-day targets. Bold indicates the minimum
within each column; it does not indicate a significance test.}
\label{tab:test}
\begin{tabular}{lrrrr}
\toprule
& \multicolumn{2}{c}{One day} & \multicolumn{2}{c}{Seven days}\\
Model & RMSE & MAE & RMSE & MAE\\
\midrule
Persistence & 2.578 & 1.978 & 5.556 & 4.371\\
Seasonal & 4.003 & 2.995 & 4.034 & 3.028\\
Nonseasonal ARIMA & 2.486 & 1.932 & 4.903 & 3.787\\
Combined & \textbf{2.405} & \textbf{1.815} & \textbf{4.019} & \textbf{3.022}\\
\bottomrule
\end{tabular}
\end{center}

Relative to persistence, the combined model reduces test RMSE by approximately
6.7\% at one day and 27.7\% at seven days. Relative to nonseasonal ARIMA,
the reductions are about 3.2\% and 18.0\%. The much larger seven-day gain over
persistence should not be attributed wholly to ARMA errors: seasonal regression
alone already produces nearly the same seven-day accuracy as the combined model.

At seven days, the combined RMSE is 4.019\,$^\circ$C and seasonal regression
is 4.034\,$^\circ$C. The unrounded difference is about 0.015\,$^\circ$C,
or 0.37\% of the seasonal baseline's RMSE. MAE differs by only about
0.005\,$^\circ$C. This is a small descriptive advantage, and validation had
ranked those two models in the opposite order. We therefore do not interpret
the held-out ordering as strong evidence that short-term dynamics materially
improve seven-day forecasting.

\begin{center}
\includegraphics[width=\linewidth]{figures/accuracy.pdf}
\captionof{figure}{Test RMSE by model and horizon. Explicit annual seasonality
is particularly valuable for seven-day forecasts; residual dynamics substantially
improve on calendar-only regression at one day.}
\end{center}

The comparisons are made on identical dates, so differences are not caused by
models receiving easier target subsets. Nevertheless, this is a single test
year at one station. We do not compute significance tests or treat overlapping
seven-day errors as independent samples. Claims concern observed forecast
accuracy under this fixed protocol, rather than universal superiority.

\newpage
\section{Prediction intervals and residual diagnostics}
\begin{center}
\captionof{table}{Combined-model nominal 95\% prediction intervals in the 2024
test. Width is measured in Celsius. Coverage is empirical over observed targets.}
\begin{tabular}{lrrr}
\toprule
Horizon & Targets & Coverage & Mean width\\
\midrule
One day & 366 & 93.7\% & 9.295\\
Seven days & 360 & 95.6\% & 14.856\\
\bottomrule
\end{tabular}
\end{center}

The seven-day intervals are wider because uncertainty accumulates as recent
residual information becomes less informative. Their average coverage is close
to nominal. One-day intervals cover only 93.7\% of targets, so improved point
accuracy does not establish complete calibration. Coverage and width must be
read together: wide intervals can achieve high coverage without being useful,
while narrow intervals may miss important temperature departures.

\begin{center}
\includegraphics[width=\linewidth]{figures/forecasts.pdf}
\captionof{figure}{Rolling test predictions and conditional 95\% intervals.
Each plotted forecast is issued at its own origin; these curves are not a
single forecast issued at the start of 2024. The seven-day forecast closely
tracks the seasonal mean and misses the sharp January cold departure.}
\end{center}

The refitted combined model has Ljung--Box p-values 0.459 and 0.158 at lags
14 and 30, respectively. These diagnostics use the longest uninterrupted
observed residual segment, 1,065 days after handling the initial burn-in and
the missing day. However, Jarque--Bera rejects normality
($p\approx5.1\times10^{-13}$), and an ARCH LM test with 14 lags gives
$p=0.025$, suggesting changing conditional variance. Thus removing much of
the serial mean dependence does not remove all distributional structure.

The reported Gaussian intervals condition on fitted coefficients and omit
parameter uncertainty. We retain these intervals to evaluate the proposed
method faithfully, rather than recalibrating them using held-out outcomes.
Distributional extensions would require a new validation protocol.

\newpage
\section{Seasonal behavior and practical limitations}
\subsection{Errors vary across target-date seasons}
We group test targets into December--February (DJF), March--May (MAM),
June--August (JJA), and September--November (SON). DJF combines January and
February 2024 with December 2024; it is a seasonal category, not a single
contiguous winter episode. There are 91, 92, 92, and 91 one-day targets in
these groups. The seven-day winter group has 85 targets because the first
six January dates are not seven-day targets under our origin rule.

\begin{center}
\includegraphics[width=\linewidth]{figures/seasons.pdf}
\captionof{figure}{Target-season RMSE. Winter is the most difficult season
at both horizons. Lines connect seasonal categories for readability and do not
represent an interpolated time series or confidence intervals.}
\end{center}

For the combined model, winter RMSE is 2.794\,$^\circ$C at one day and
5.616\,$^\circ$C at seven days. The seven-day RMSE is 3.534 in spring,
3.579 in summer, and 2.977 in autumn. The full-year score therefore hides
substantial seasonal variation. The sharp January cold departure visible in
the forecast plot contributes to the winter difficulty, but this descriptive
analysis does not establish a causal explanation for every error.

Winter combined-model interval coverage is 91.2\% at one day and 87.1\%
at seven days, considerably below the nominal 95\%. Seven-day coverage in
spring, summer, and autumn is 98.9\%, 96.7\%, and 98.9\%, respectively.
Near-nominal aggregate coverage consequently combines winter undercoverage
with relatively high coverage elsewhere. For outdoor operations, aggregate
calibration alone would be insufficient to assess risk during cold episodes.

\subsection{What this experiment can and cannot establish}
The forecast uses station history and calendar information only. It cannot
anticipate a changing atmospheric regime from meteorological predictors because
none are supplied. This makes it a transparent statistical benchmark, but not
a replacement for a comprehensive operational weather forecast. We do not
compare against numerical weather prediction or claim operational superiority.

The data scope is six years, with four years in the initial fit and one year
each for selection and testing. A single station limits geographic
generalization. Station changes, reporting changes, or a different run of weather
could affect performance; we audit coverage but do not independently validate
every source measurement. The fixed annual curve also does not model a secular
climate trend or year-specific seasonal shifts.

The frozen-parameter test isolates state adaptation from re-estimation, but it
does not show whether periodic parameter refitting would help. Our results
also assume completed daily observations are available immediately at the
origin. A deployed system would need to align origins with actual publication
times. These constraints should accompany any use of the forecasts for planning.

\newpage
\section{Reproducibility and conclusion}
\subsection{Artifacts and executable checks}
The project includes scripts for the station audit, validation search, frozen
test, and result summaries, plus an executed demonstration notebook.
Raw NOAA station-year files are cached locally. Download manifests record URLs
and checksums; reproduction hashes identify the processed data, analysis code,
protocol, and metric table. Saved candidate forecasts and JSON parameter records
make the selected results inspectable without rerunning the entire search.
The recorded analysis environment uses Python 3.13.12, NumPy 2.4.3,
pandas 3.0.1, statsmodels 0.14.6, and Matplotlib 3.10.8.

The notebook reconstructs a historical forecast using frozen parameters and
only observations through a selected origin, then checks both point forecasts
and interval endpoints against the saved rolling evaluation. Protocol tests
verify consistent Fourier phase across a leap day, exact horizon offsets,
unchanged fitted parameters, and insensitivity of earlier forecasts to changed
future observations. All ten notebook code cells execute successfully. These
checks verify implementation properties; they do not establish generalization
beyond the evaluated dataset. Code is prepared locally; public GitHub hosting
remains a separate deliverable.

\subsection{Conclusion}
Combining annual Fourier regression with ARMA(1,1) errors yields the lowest
observed one-day validation and test RMSE among the compared models.
At seven days, explicit annual seasonality provides nearly all of the advantage
over persistence and nonseasonal ARIMA; the combined model's small test gain
over seasonal regression should be interpreted cautiously. Thus the answer to
the project question depends on horizon rather than supporting an unqualified
benefit from additional dynamics.

Prediction intervals reveal a second distinction. Good average performance
does not guarantee reliable uncertainty during difficult periods: winter
coverage is poor even when annual seven-day coverage is near 95\%. Residual
nonnormality and changing variance provide diagnostic reasons to investigate
more flexible uncertainty models in future work. Such extensions, additional
stations, and repeated future-year evaluation would require new validation
rather than tuning to the completed 2024 test.

\section*{References}
\raggedright
\begin{enumerate}
\item NOAA National Centers for Environmental Information. \emph{Global Surface
Summary of the Day (GSOD)}. Dataset and field documentation.
\url{https://www.ncei.noaa.gov/data/global-summary-of-the-day/access/};
\url{https://www.ncei.noaa.gov/data/global-summary-of-the-day/doc/readme.txt}.
\item R. J. Hyndman and G. Athanasopoulos. \emph{Forecasting: Principles and
Practice}, 3rd ed. OTexts, 2021. Chapter 10, Dynamic regression models,
including dynamic harmonic regression.
\url{https://otexts.com/fpp3/dynamic-regression.html}.
\item The statsmodels Developers. \emph{SARIMAX: State-Space Time-Series
Modeling}, version 0.14.6. Model and forecast-state update documentation.
\url{https://www.statsmodels.org/v0.14.6/generated/statsmodels.tsa.statespace.sarimax.SARIMAX.html}.
\item R. J. Hyndman and Y. Khandakar. Automatic time series forecasting:
the forecast package for R. \emph{Journal of Statistical Software}, 27(3):1--22,
2008. \url{https://doi.org/10.18637/jss.v027.i03}.
\item S. J. Taylor and B. Letham. Forecasting at scale.
\emph{The American Statistician}, 72(1):37--45, 2018.
\url{https://doi.org/10.1080/00031305.2017.1380080}.
\end{enumerate}
\end{document}
